% Part: set-theory
% Chapter: z
% Section: pairs

\documentclass[../../../include/open-logic-section]{subfiles}

\begin{document}

\olfileid{sth}{z}{pairs}
\olsection{যুগল}

পরের স্বতঃসিদ্ধটি এই:

\begin{axiom}[যুগল]
যেকোনো সেট $a, b$-এর জন্য $\{a, b\}$ সেটটি আছে।
\[
  \forall a \forall b \exists P \forall x (x \in P \liff (x = a \lor x = b))
\]
\end{axiom}

পর্যায়ক্রমিক ধারণা দিয়ে এই স্বতঃসিদ্ধটির যথার্থতা
এভাবে বোঝা যায়। ধরি, $a$ পর্যায়~$S$-এ এবং $b$
পর্যায়~$T$-এ উপলব্ধ। $S$ ও~$T$-এর মধ্যে যে পর্যায়টি
পরে আসে, সেটিকে~$M$ বলি। তাহলে $a$ ও~$b$ উভয়ই
পর্যায়~$M$-এ উপলব্ধ; তাই $M$-এর পরের যেকোনো
পর্যায়ে $\{a,b\}$ একটি গঠনযোগ্য সংগ্রহ।

কিন্তু একটু থামো!{} $M$-এর পরে আদৌ কোনো পর্যায়
\emph{আছে} বলে ধরব কেন? না থাকলে আমাদের যুক্তি
ব্যর্থ হবে। তাই যুগলের স্বতঃসিদ্ধকে সমর্থন করতে
\olref[sth][z][story]{sec}-এর কাহিনিতে আরেকটি নীতি
যোগ করতে হবে:
\begin{enumerate}
  \item[] \stagessucc. কোনো শেষ পর্যায় নেই।
\end{enumerate}
এই নীতির ভিত্তি কী? শোয়েনফিল্ডের বিবরণে
\emph{স্পষ্টভাবে} বলা হয়নি যে শেষ পর্যায় নেই।
কঠোর অর্থে এটি কাহিনিতে একটি অতিরিক্ত সংযোজন
হলেও পর্যায়ে পর্যায়ে সেট গঠনের মৌলিক ধারণার
সঙ্গে ভালোভাবেই মেলে। পরের আলোচনায় নীতিটি
আমরা গ্রহণ করব। তাই যুগলের স্বতঃসিদ্ধও গ্রহণ করব।

নতুন স্বতঃসিদ্ধটি হাতে থাকলে আরও অনেক সেটের
অস্তিত্ব প্রমাণ করা যায়। যেমন:

\begin{prop}\ollabel{prop:pairsconsequences}
যেকোনো সেট $a$ ও~$b$-এর জন্য নিচের সেটগুলি আছে:
  \begin{enumerate}
    \item\ollabel{singleton} $\{a\}$
    \item\ollabel{binunion} $a \cup b$
    \item\ollabel{tuples} $\tuple{a, b}$
  \end{enumerate}
\end{prop}

\begin{proof}
\olref{singleton}. যুগলের স্বতঃসিদ্ধে $\{a, a\}$ আছে;
সদস্যভিত্তিক সমতা অনুযায়ী সেটিই~$\{a\}$।

\olref{binunion}. যুগলের স্বতঃসিদ্ধে $\{a, b\}$
আছে। এবার সংযোগের স্বতঃসিদ্ধে
$a \cup b = \bigcup \{a, b\}$ আছে।

\olref{tuples}. \olref{singleton} অনুযায়ী $\{a\}$
আছে। যুগলের স্বতঃসিদ্ধে $\{a, b\}$ আছে। এবার
আবার যুগলের স্বতঃসিদ্ধে
$\{\{a\}, \{a, b\}\} = \tuple{a, b}$ আছে।
\end{proof}

\begin{prob}
দেখাও যে, যেকোনো সেট $a, b, c$-এর জন্য
$\{a, b, c\}$ সেটটি আছে।
\end{prob}

\begin{prob}
দেখাও যে, যেকোনো সেট $a_1, \ldots, a_n$-এর
জন্য $\{a_1, \ldots, a_n\}$ সেটটি আছে।
\end{prob}

%\begin{proof} By two applications of Pairs, $\{\{a_1, a_2\}, \{a_1,
%   a_3\}\}$ exists. By Union and Extensionality, $\bigcup \{\{a_1,
%   a_2\}, \{a_1, a_3\}\} = \{a_1, a_2, a_3\}$ exists. Repeat this
%   trick as often as necessary. \end{proof}

\end{document}
